% Copyright (c) 2026 Eugene Joseph M. Ragasa.
% SPDX-License-Identifier: Apache-2.0
% Author: Eugene Joseph M. Ragasa
% Formal proof completed: 2026-10-27
\documentclass[11pt]{article}
\usepackage[T1]{fontenc}
\usepackage{amsmath,amssymb}
\usepackage[margin=1in]{geometry}
\usepackage[hidelinks]{hyperref}
\title{Proof Dossier: Sewing Closure of Operator-Field Differences}
\author{Eugene Joseph M. Ragasa}
\date{Formal proof completed: 2026-10-27}
\begin{document}
\raggedright
\maketitle

\section{Identity and status}
Statement ID \texttt{BG-BLOCH-THM-013}; status \texttt{proved; no admissions}. Lean declaration
\path{ProjectKoios.BandGauge.Bloch.isBlochOperatorField_operatorFieldDifference} in
\path{ProjectKoios/BandGauge/Bloch/AlignedDifference.lean}.

\section{Exact mathematical statement}
If two operator fields obey the same sewing covariance, then their pointwise
reference-minus-candidate difference obeys that sewing covariance as well.

\section{Assumptions and conventions}
One group action, one sewing transport, and two operator fields on the same complex fiber are
supplied. Each field has its own explicit proof of covariance with that same sewing data.

\section{Mathematical argument}
Evaluate the difference field at the transformed point and a transported vector. Apply the
reference sewing law and the candidate sewing law separately. The result is the difference of
two transported vectors. Linearity of the sewing transport converts that difference into the
transport of the original pointwise difference, which is precisely the required covariance law.

\section{Lean realization}
The proof introduces the group element, point, and vector, unfolds the target to the two
pointwise differences, rewrites with both sewing hypotheses, and closes with
\texttt{map\_sub}. It was checked with Lean \texttt{4.35.0} and Mathlib commit
\texttt{56a005855ea17e0b0958f73ec801ae140b0a1679}.

\section{Provenance and contribution boundary}
This is source-neutral finite covariance algebra at Formal merge
\texttt{b998c09550161193e8edbe3aeaa9d93d183693fa}. No external manuscript controls the Lean
statement, and no priority or novelty claim is made.

\section{Audit result}
The named axiom report lists \texttt{propext}, \texttt{Classical.choice}, and
\texttt{Quot.sound} and no repository-owned axiom. The repository admission audit found no
\texttt{sorry} or \texttt{admit} in the checked source.

\section{Boundary and nonclaims}
Sewing closure does not prove continuity, smoothness, self-adjointness, numerical conformity,
optimizer correctness, topology, or physical adequacy.

\section{Review status}
AI-assisted exposition, checked against the Lean source and theorem manifest. Kernel checking
establishes the exact proposition only. No independent expert mathematical, scientific,
novelty, or publication review is recorded by this dossier.

\end{document}
